\PassOptionsToPackage{unicode}{hyperref}
\PassOptionsToPackage{hyphens}{url}
\documentclass[11pt,a4paper]{article}

\usepackage[T1]{fontenc}
\usepackage[utf8]{inputenc}
\usepackage[brazil]{babel}
\usepackage{lmodern}
\usepackage{microtype}

\usepackage{xcolor}
\definecolor{javaOrange}{HTML}{E76F00}
\definecolor{javaDeep}{HTML}{BF5700}
\definecolor{javaBlue}{HTML}{1565C0}
\definecolor{javaTeal}{HTML}{00695C}
\definecolor{javaAmber}{HTML}{B45309}
\definecolor{javaInk}{HTML}{2B2140}
\definecolor{javaCodeBg}{HTML}{FFF8F0}
\definecolor{javaRule}{HTML}{FFD7B5}
\definecolor{javaShade}{HTML}{FFF3E8}

\usepackage{tcolorbox}
\tcbuselibrary{skins,breakable,listings}
\usepackage{listings}
\usepackage{fvextra}

\lstdefinestyle{javastyle}{
  language=Java,
  basicstyle=\ttfamily\small,
  keywordstyle=\color{javaBlue}\bfseries,
  stringstyle=\color{javaDeep},
  commentstyle=\color{javaTeal}\itshape,
  breaklines=true,
  showstringspaces=false,
  tabsize=4,
}

\newtcblisting{javacode}{
  listing only,
  listing style=javastyle,
  breakable, enhanced,
  colback=javaCodeBg, colframe=javaDeep,
  leftrule=4pt, rightrule=0.4pt, toprule=0.4pt, bottomrule=0.4pt,
  arc=2pt, fontupper=\small,
  top=4pt, bottom=4pt, left=6pt, right=6pt,
}

\newtcbox{\inlinecode}{
  on line,
  colback=javaCodeBg, colframe=javaRule,
  boxrule=0.5pt, arc=2pt,
  fontupper=\ttfamily\small\color{javaDeep},
  boxsep=1pt, left=2pt, right=2pt, top=1pt, bottom=1pt,
}

\usepackage[margin=2.4cm, top=2.8cm, bottom=2.8cm]{geometry}

\usepackage{fancyhdr}
\pagestyle{fancy}
\fancyhf{}
\renewcommand{\headrulewidth}{0.4pt}
\renewcommand{\headrule}{\hbox to\headwidth{\color{javaRule}\leaders\hrule height \headrulewidth\hfill}}
\fancyhead[L]{\small\sffamily\color{javaDeep} Java Programming MOOC}
\fancyhead[R]{\small\sffamily\color{javaDeep} Lição 2}
\fancyfoot[C]{\small\sffamily\color{gray}\thepage}

\usepackage{titlesec}
\titleformat{\section}
  {\sffamily\Large\bfseries\color{javaOrange}}{\thesection}{0.7em}{}
  [\vspace{2pt}{\color{javaRule}\titlerule[1.2pt]}]
\titleformat{\subsection}
  {\sffamily\large\bfseries\color{javaDeep}}{\thesubsection}{0.6em}{}
\titleformat{\subsubsection}
  {\sffamily\normalsize\bfseries\color{javaTeal}}{\thesubsubsection}{0.5em}{}

\usepackage{enumitem}
\setlist[itemize]{itemsep=2pt, topsep=4pt, parsep=0pt}
\setlist[enumerate]{itemsep=2pt, topsep=4pt, parsep=0pt}

\usepackage{booktabs}
\usepackage{array}
\usepackage{colortbl}

\usepackage{hyperref}
\hypersetup{
  colorlinks=true,
  linkcolor=javaBlue, urlcolor=javaBlue, citecolor=javaBlue,
  pdftitle={Programação em Java — Lição 2: Repetição e Métodos},
  pdfauthor={Java Programming MOOC · Universidade de Helsinque},
  pdflang={pt-BR},
}

\newtcolorbox{destaque}[1][Atenção]{
  enhanced, breakable,
  colback=javaShade, colframe=javaOrange,
  leftrule=4pt, rightrule=0.4pt, toprule=0.4pt, bottomrule=0.4pt,
  arc=2pt,
  fonttitle=\sffamily\bfseries\color{javaOrange},
  title=#1,
  top=4pt, bottom=4pt,
}

\makeatletter
\newcommand{\subtitle}[1]{\gdef\@subtitle{#1}}
\renewcommand{\maketitle}{%
  \begin{center}
    {\color{javaRule}\rule{\linewidth}{1pt}}\\[6pt]
    {\Huge\sffamily\bfseries\color{javaOrange} \@title}\\[6pt]
    \ifx\@subtitle\undefined\else
      {\Large\sffamily\color{javaDeep} \@subtitle}\\[4pt]
    \fi
    {\small\sffamily\color{gray} \@author}\\[2pt]
    {\small\sffamily\color{gray} \@date}\\[6pt]
    {\color{javaRule}\rule{\linewidth}{1pt}}
  \end{center}
  \vspace{1em}
}
\makeatother

\title{Programação em Java}
\subtitle{Lição 2: Repetição e Métodos}
\author{Java Programming MOOC · Universidade de Helsinque — tradução para o português}
\date{2026}

\begin{document}
\maketitle
\tableofcontents
\vspace{1.5em}

% ─────────────────────────────────────────────────────────────────────────────
\section{Objetivos de aprendizagem}

Ao final desta lição, você será capaz de:

\begin{itemize}
  \item Usar o laço \inlinecode{while} para repetir código
  \item Usar os comandos \inlinecode{break} e \inlinecode{continue}
  \item Usar o laço \inlinecode{for} para percorrer sequências
  \item Criar e chamar métodos próprios
  \item Criar métodos com parâmetros e com valor de retorno
\end{itemize}

% ─────────────────────────────────────────────────────────────────────────────
\section{Padrões comuns de repetição}

Antes de aprender os laços, é útil reconhecer padrões frequentes.

\textbf{Leitura contínua até um valor sentinela:}

\begin{javacode}
Scanner scanner = new Scanner(System.in);
while (true) {
    String linha = scanner.nextLine();
    if (linha.equals("fim")) {
        break;
    }
    System.out.println("Você digitou: " + linha);
}
\end{javacode}

\textbf{Acúmulo de resultado:}

\begin{javacode}
int soma = 0;
int contador = 0;
// ... leitura de valores no laço ...
double media = (double) soma / contador;
\end{javacode}

% ─────────────────────────────────────────────────────────────────────────────
\section{O laço \texttt{while}}

O laço \inlinecode{while} repete um bloco de código enquanto uma condição for verdadeira.

\subsection{Forma básica}

\begin{javacode}
int numero = 1;
while (numero <= 5) {
    System.out.println(numero);
    numero++;
}
\end{javacode}

\subsection{Laço infinito com \texttt{break}}

\begin{javacode}
Scanner scanner = new Scanner(System.in);
int soma = 0;
while (true) {
    System.out.print("Digite um número (0 para sair): ");
    int n = Integer.valueOf(scanner.nextLine());
    if (n == 0) {
        break;
    }
    soma += n;
}
System.out.println("Soma total: " + soma);
\end{javacode}

\subsection{O comando \texttt{continue}}

O \inlinecode{continue} pula o restante da iteração atual e volta ao início do laço:

\begin{javacode}
int i = 0;
while (i < 10) {
    i++;
    if (i % 2 == 0) {
        continue; // pula os pares
    }
    System.out.println(i); // imprime apenas ímpares
}
\end{javacode}

\subsection{Calculando estatísticas}

Variáveis de acúmulo devem ser declaradas \textbf{antes} do laço:

\begin{javacode}
Scanner scanner = new Scanner(System.in);
int soma = 0, quantidade = 0, maior = Integer.MIN_VALUE;

while (true) {
    System.out.print("Digite um número (ou -1 para encerrar): ");
    int n = Integer.valueOf(scanner.nextLine());
    if (n == -1) break;
    soma += n;
    quantidade++;
    if (n > maior) maior = n;
}

if (quantidade > 0) {
    System.out.println("Quantidade: " + quantidade);
    System.out.println("Soma: " + soma);
    System.out.println("Média: " + (double) soma / quantidade);
    System.out.println("Maior: " + maior);
}
\end{javacode}

% ─────────────────────────────────────────────────────────────────────────────
\section{O laço \texttt{for}}

O \inlinecode{for} é mais compacto e muito usado quando se sabe antecipadamente quantas iterações serão feitas.

\subsection{Sintaxe}

\begin{javacode}
for (inicialização; condição; atualização) {
    // corpo
}
\end{javacode}

\subsection{Exemplos}

\begin{javacode}
// Conta de 1 a 5
for (int i = 1; i <= 5; i++) {
    System.out.println(i);
}

// Pulos de 2 em 2
for (int i = 0; i <= 10; i += 2) {
    System.out.println(i); // 0 2 4 6 8 10
}
\end{javacode}

\subsection{Laços aninhados}

\begin{javacode}
for (int linha = 1; linha <= 3; linha++) {
    for (int coluna = 1; coluna <= 3; coluna++) {
        System.out.print(linha + "x" + coluna + "=" + (linha*coluna) + "  ");
    }
    System.out.println();
}
\end{javacode}

% ─────────────────────────────────────────────────────────────────────────────
\section{Métodos}

Um \textbf{método} é um bloco de código nomeado que realiza uma tarefa específica.

\subsection{Método sem parâmetros e sem retorno}

\begin{javacode}
public static void saudacao() {
    System.out.println("Olá! Bem-vindo ao programa.");
}
// chamada:
saudacao();
\end{javacode}

\subsection{Método com parâmetros}

\begin{javacode}
public static void saudar(int vezes) {
    for (int i = 0; i < vezes; i++) {
        System.out.println("Olá!");
    }
}
\end{javacode}

\begin{destaque}[Atenção]
Parâmetros são \textbf{cópias} dos valores passados. Modificá-los dentro do método não altera as variáveis originais.
\end{destaque}

\subsection{Método com valor de retorno}

\begin{javacode}
public static int soma(int a, int b) {
    return a + b;
}

public static double media(int a, int b, int c) {
    return (a + b + c) / 3.0;
}

public static boolean ehPar(int n) {
    return n % 2 == 0;
}
\end{javacode}

\subsection{Exemplo completo}

\begin{javacode}
import java.util.Scanner;

public class Calculadora {

    public static double calcularMedia(int a, int b, int c) {
        return (a + b + c) / 3.0;
    }

    public static int calcularMaximo(int a, int b) {
        return a > b ? a : b;
    }

    public static void main(String[] args) {
        Scanner scanner = new Scanner(System.in);
        System.out.print("Primeiro número: ");
        int n1 = Integer.valueOf(scanner.nextLine());
        System.out.print("Segundo número: ");
        int n2 = Integer.valueOf(scanner.nextLine());
        System.out.print("Terceiro número: ");
        int n3 = Integer.valueOf(scanner.nextLine());

        System.out.println("Média: " + calcularMedia(n1, n2, n3));
        System.out.println("Maior dos dois primeiros: " + calcularMaximo(n1, n2));
    }
}
\end{javacode}

% ─────────────────────────────────────────────────────────────────────────────
\section{Quando usar \texttt{while} vs.\ \texttt{for}}

\begin{center}
\begin{tabular}{l l}
\toprule
\textbf{Situação} & \textbf{Use} \\
\midrule
Número de iterações desconhecido & \texttt{while} \\
Lendo até um valor sentinela & \texttt{while (true)} com \texttt{break} \\
Percorrendo um intervalo conhecido & \texttt{for} \\
Percorrendo uma coleção & \texttt{for-each} \\
\bottomrule
\end{tabular}
\end{center}

% ─────────────────────────────────────────────────────────────────────────────
\section{Resumo}

\begin{itemize}
  \item O laço \inlinecode{while} repete enquanto a condição for verdadeira; \inlinecode{break} sai, \inlinecode{continue} pula.
  \item O laço \inlinecode{for} é ideal quando sabemos o número de iterações.
  \item Variáveis de acúmulo devem ser declaradas \textbf{antes} do laço.
  \item Um \textbf{método} pode receber parâmetros e retornar valores.
  \item Parâmetros são cópias: alterá-los no método não afeta as variáveis originais.
  \item \inlinecode{return} encerra o método e devolve um valor.
\end{itemize}

% ─────────────────────────────────────────────────────────────────────────────
\section{Exercícios}

\begin{enumerate}
  \item \textbf{Tabuada} — Lê um número inteiro \(n\) e imprime a tabuada de \(n\) (de 1 a 10).

  \item \textbf{Fatorial} — Crie um método \inlinecode{fatorial(int n)} que retorna o fatorial de \(n\). Teste com \(0, 1, 5\) e \(10\).

  \item \textbf{Contagem de pares e ímpares} — Lê números inteiros até o usuário digitar \(0\) e exibe quantos foram pares e quantos foram ímpares.

  \item \textbf{Potência} — Crie \inlinecode{potencia(int base, int expoente)} usando um laço (sem \inlinecode{Math.pow}). Teste com \inlinecode{potencia(2, 10)}.

  \item \textbf{Validação de entrada} — Crie \inlinecode{lerInteiroPositivo(Scanner scanner)} que pede um número positivo, repete se inválido, e retorna o número válido. Use-o para ler dois números e exibir o produto.

  \item \textbf{Estrelas} — Crie \inlinecode{imprimirEstrelas(int n)} que imprime \(n\) asteriscos. Use-o para imprimir um triângulo crescente de 1 a 5 linhas.
\end{enumerate}

\end{document}
